\section*{Introduction}
Model-based phylogenetic inference has become routine, largely due to the availability of phylogenetic software like PHYLIP \citepMain{felsenstein2005}, PhyML \citepMain{guindon2010}, RAxML \citepMain{Stamatakis2014} or IQ-Tree \citepMain{minh2020}. This process involves inferring a phylogenetic tree from an alignment and subsequently validating it using various statistical methods. For example \citepMain{Anisimova2006} provide a likelihood ratio test to evaluate the reliability of short branches. The bootstrap by \citepMain{felsenstein1985confidence} offers a support value for a branch in the tree, where a low bootstrap value indicates that the branch is not well support. Typically, branches with high-bootstrap support are considered as a real phylogenetic signal, justifying the separation between the two clusters in the tree. Implicit to all these analyses is the assumption that the alignment contains sufficient phylogenomic information to conclude that the branch displays a historical relationship.

\medskip
However, the antagonist to phylogenetic information is saturation, which occurs when there are too many substitutions in an alignment to provide reliable information about true genetic distances \citepMain{philippe2011}. For two sequences, an alignment is saturated if we cannot distinguish if the two sequences are homologous or were just sampled at random independently  \citepMain{jermiin2020, philippe1999, brown1982}. When two sequences are saturated, connecting them by a branch to display the smallest possible tree lacks phylogenetic meaning. 

\medskip
For alignments with more than two taxa and a tree with many branches, it becomes crucial to identify which branches are subject to saturation and, more importantly, how to measure this saturation.  If we detect saturation or loss of phylogenetic information for a branch,  then the phylogenetic information for the two clusters/splits generated by the branch has eroded.

\medskip

Leaving aside a maximum-parsimony test \citepMain{Archie1989} or ad hoc analysis \citepMain{liu2014, davalos2008}, current techniques for checking saturation \citepMain{duchene2021, xia2003index} study entire alignments and their phylogenetic content. These methods cannot test for saturation of individual branches in a phylogenetic tree.

\medskip 
Here we introduce SatuTe (SATUration TEst), which addresses this gap. For a given tree and alignment, SatuTe evaluates whether the alignments provide enough phylogenetic information to justify the branches present in the tree. SatuTe embeds the analysis of saturation between two sequences into a framework that handles multiple sequences and a tree. Thus, this approach opens a new avenue for studying the relationship between alignments and  trees.

\medskip 

In the following, we present the theoretical foundations of SatuTe. We explain the test designed to reject the null hypothesis of no phylogenetic signal/saturation in favor of the alternative hypothesis of phylogenetic signal present/no saturation. Then, we use simulated data to evaluate the performance of the test and its applicability across various phylogenetic application scenarios. These scenarios include cases where a given  tree and some alignment  are analyzed to determine which branches are supported by a phylogenetic signal from the alignment, with branch lengths provided or inferred from the alignment using maximum likelihood methods. Lastly, we apply the test to real data by analyzing the Tree of Life (ToL) for ribosomal protein and RNA alignments. Focusing on the branch leading to Eukaryota, our proposed analyses contributes to the two-domain hypothesis \citepMain{williams2D, Hug2016} and contrasts it with the "classical" ribosomal RNA-based tree \citepMain{woese, gribaldo2010}, which supports the three-domain hypothesis.



%Then, we discuss the performance of the test for three typical questions:
%\begin{enumerate}
%    \item[(i)] \textbf{A given tree and alignment:} We ask which branches in the alignment are supported by a phylogenetic signal. The branch lengths are either given or inferred from the alignment using maximum likelihood.
%    \item[(ii)] \textbf{A maximum likelihood tree inferred from the alignment:} We ask which regions (sliding windows) are phylogenetic informative for the branches in the tree.
%    \item[(iii)] \textbf{A similar scenario as (ii):} We ask for which branches in the tree the alignment is informative. Contrary to (i) and (ii), a Bonferroni test for multiple testing is required.
%\end{enumerate}
%
%\medskip
%Finally, we exemplarily analyze the branch connecting Eukaryotes with the rest of the tree of life (ToL) for ribosomal proteins, which led to the two-domains hypothesis \citepMain{williams2D, Hug2016}, and the “classical” ribosomal RNA-based tree \citepMain{woese,gribaldo2010}, which led to the three-domains hypothesis.


\section*{Results}
\subsection*{Conceptual basis of SatuTe}\label{sec:SatuTe}

Consider an alignment of length $n$, where the columns represent $n$ (site) patterns $\partial$, and each of the $m$ rows corresponds to a DNA or amino acid sequence of the taxa (Fig.\,\ref{fig:panel1}a), which form the leaf-set of a phylogenetic tree $\mathbb{T}$.

We assume a reversible and stationary model of sequence evolution \citepMain{tavare1986} with probability transition matrix $\bm{P}(t)$ and stationary distribution $\bm{\pi}$. The matrix $\bm{P}(t) = e^{\bm{Q}\cdot t}$ with rate matrix $\bm{Q}$ describes the probability of a nucleotide or amino acid pair after $t$ substitutions per site have occurred. 

\medskip

Given the alignment (Fig.\,\ref{fig:panel1}a), the model of sequence evolution $\bm{Q}$, and a phylogenetic tree $\mathbb{T}$, we examine whether the branch $AB$ (Fig.\,\ref{fig:panel1}b) is supported by the phylogenetic information derived from the topology and subalignments of the branch-induced subtrees $\mathbb{T}_A$ and $\mathbb{T}_B$.

\medskip

To this end, using the Pruning algorithm \citepMain{Felsenstein1981}, we compute the (partial) likelihood vectors $\va$ and $\vb$  at nodes A and B (Fig.\,\ref{fig:panel1}b) for each site pattern $\partial$. These vectors represent the partial likelihoods for the sub-patterns  $\va$ and $\vb$  as induced by the branch AB. The probability of observing the full alignment pattern $\partial$ is then computed as follows:

\begin{equation} \label{eq:branch_likelihood_given_t} \Prob(\partial \mid t)  = \va^T \cdot\diag (\bm{\pi})\cdot \bm{P}(t) \cdot\vb.
\end{equation}


In addition, the model of sequence $\bm{Q}$ evolution allows for a spectral decomposition \citepMain{levin2017markov}, such that

\begin{equation} \label{eq:spectral_decomposition}
   \bm{P}(t) = \bm{1}\cdot\bm{\pi}^T + \sum_{i=1}^3\bm{v_i}\cdot \bm{h_i}^T e^{\lambda_it},
\end{equation}
 
where  ${\lambda_0 = 0 > \lambda_1 \geq \lambda_2 \geq \lambda_3}$  are the eigenvalues of the rate matrix $\bm{Q}$ and $\bm{h_k}$,$\bm{v_k}$  are the left and right eigenvectors associated with eigenvalue $\lambda_k$, respectively. 

After plugging \eqref{eq:spectral_decomposition} into \eqref{eq:branch_likelihood_given_t} and some formula manipulations (see online methods \ref{msec:branch likelihood}), we obtain
\begin{equation} \label{eq:branch_likelihood_at_site_t} \frac{\Prob(\partial \mid t)}{\Prob(\pred)\Prob(\pblue)} =  1 + \sum_{i =1}^3 C^\partial_i\cdot e^{\lambda_i t}.
\end{equation}


This reformulation reveals that a phylogenetic tree with an increasingly long branch separating two subtrees $\mathbb{T}_A$ and $\mathbb{T}_B$   produces the same pattern probabilities as a disconnected tree where subtrees $\mathbb{T}_A$ and $\mathbb{T}_B$ are independent. Thus, our null hypothesis is that the alignment was generated by independent subtrees (branch saturation). 

\medskip
We call $C^\partial_i$ the coherence coefficient (for the non-zero eigenvalues $i$). In the online methods section \ref{msec:branch likelihood}, we demonstrate that the expectation of the coherence coefficient equals zero if the subalignments were independently evolving in the subtrees. This implies that there is no phylogenetic information between the subalignments, and  the branch as indication of a common evolutionary history is not justified.

We also observe that the contribution of the coherence coefficient to the sum decreases exponentially fast. Therefore, the coherence coefficient $C^\partial_1$ contributes most significantly to the sum, and we will restrict our discussion to it.

\medskip 
By averaging the coherence values $C^\partial_1$  over the alignment lengths and computing the standard deviation (see online methods \ref{msec:test}), we obtain the average coherence coefficient $\hat{C}_1$. According to the central limit theorem, this average coherence coefficient is normally distributed (Fig.\,\ref{fig:panel1}c) and can be transformed into a z-score. If the z-score is larger than the critical value $z_\alpha$ for a specified significance level $\alpha$, the null hypothesis of independently evolving subalignments is rejected. This indicates that the subalignments share enough phylogenetic information to justify the connection in the tree (Fig.\,\ref{fig:panel1}d).


We note that the calculation of the mean coherence does not depend on the edge length with which the two subtrees are connected. Here the similarity of the test to the test of saturation between two sequences becomes clear, where the counted nucleotide or amino acid pairs decide on saturation. Our test is therefore directly applicable to two sequences, this is obvious when considering Eq \eqref{eq:branch_likelihood_given_t}. But Eq \eqref{eq:branch_likelihood_given_t}  describes the general case.

\medskip

	\begin{figure}[H]
            \centering
            \includegraphics[width=\linewidth]{figure_1_2024_06_12.pdf}
		     \caption{
		   \textbf{The theoretical foundations of SatuTe.} \textbf{a}, In an alignment of the sequences from five taxa, each column represents a pattern $\partial$. \textbf{b}, The branch $AB$ splits the five-taxon tree into subtrees  $\mathbb{T}_A$ with sequences $S1,S4,S5$ and subtree $\mathbb{T}_B$ with sequences $S2, S3$. Additionally, branch $AB$ splits each pattern, such as $\partial = \text{CTTTC}$, into sub-patterns ${\partial A=\text{CTC}}$ and ${\partial B=\text{TT}}$. These sub-patterns are used to compute the likelihood vectors $\va$ and $\vb$, which are then utilized to compute the coherence coefficients. \textbf{c}, The average coherence $\hat{C}_1$ is approximately normal distributed, with expectation zero and standard deviation $\hat{\sigma}_1/\sqrt{n}$, if the subtrees  $\mathbb{T}_A$ and  $\mathbb{T}_B$ are evolving independently. The z-transform of the average coherence is then used to test for phylogenetic information. \textbf{d}, The null hypothesis of independently evolving subtrees is rejected at the significance level $\alpha$, if $Z>z_\alpha$. Rejection implies that the subtrees are not evolving independently and can be connected by a branch. Consequently, the subalignments, along with the subtrees, provide sufficient phylogenetic information.}
		     \label{fig:panel1}
        \end{figure}



\subsection*{SatuTe applied to simulated data}

To evaluate the performance of test for phylogenetic information using SatuTe, we analyzed a wide range of parameter selections. These included different numbers of taxa (five and sixteen), varying alignment lengths ($100$, $1000$, $10000$), different branch types to test (internal and external), and different branch lengths (see online methods \ref{msec:figprotocols}).

For the different simulation trees, we focused on the branch $AB$ and varied the length of this branch. For each parameter combination, we generated 1000 DNA alignments of a specific length along the specified simulation tree under the JC model. Using these datasets, we analyzed various phylogenetic application scenarios. Fig.\,\ref{fig:simulation} summarizes the main findings.

\medskip

In the \textbf{first scenario}, representing the ideal situation (true-tree, {\color{truetree} light red}, Fig.\,\ref{fig:simulation}), an alignment is simulated along the true (simulation) tree with its specified branch lengths. Subsequently, the phylogenetic information of the simulated alignment is tested for the external edge (Fig.\,\ref{fig:simulation}a) and the internal edge (Fig.\,\ref{fig:simulation}b). Regardless of the specific parameters of the simulation, the fraction of informative alignments for the branch $AB$ decreases from 100\% to the significance level of 5\% as the branch length increases. Thus, long branch lengths in the true evolutionary process lead to a reduction in rejection rates, or equivalently, to more instances where the subtrees do not convey phylogenetic information. According to standard statistical theory, longer alignments preserve phylogenetic information for longer branches. Therefore, the overall behavior of the curves follows a typical power curve and serves as a basis for other scenarios that deviate from this gold standard.

\medskip 

In the \textbf{second scenario}, a more realistic approach (true-tree with ML-estimated branch lengths, {\color{truetreeML} dark red}, Fig.\,\ref{fig:simulation}), the true-tree topology is known, but the branch lengths are estimated using maximum likelihood.  As shown in Fig.\,\ref{fig:simulation}, the dark red power curves are extremely close to the gold standard (true-tree scenario). For short sequences (100bp), the test appears to be a bit too liberal, with a simulated rejection level of approximately 10\%, but this is still considered tolerable.

A typical question here would be to check which branches in a generally accepted tree are supported by a new alignment for a sequence family that has not yet been investigated. In other words, where in the accepted tree topology is the new alignment phylogenetic informative?

Scenario two also allows for another even more interesting application of SatuTe. After reconstructing a phylogenetic tree from a very long phylogenomic alignment, one might be interested in identifying which genes or other structural subsequences of the alignment are informative for the branches of the reconstructed tree. Scenario two demonstrates that this approach is valid even if the respective genes constitute only a small percentage of the total alignment.

\medskip
The \textbf{third scenario}  (ML-tree, {\color{MLtree} light green}, Fig.\,\ref{fig:simulation}) involves reconstructing a maximum likelihood (ML) tree using an alignment and then determining where in the tree the alignment is informative. This scenario is the most complex due to the varying influences of ML-inference on our test. Considering all possible placements of nodes  $A$ and $B$ in $\mathbb{T}_A$ and $\mathbb{T}_B$ respectively, the probability of finding random similarities increases with the number of taxa. Another factor contributing to this high rejection rate may be the systematic underestimation of branch lengths to maximize likelihood, as shown in Supplementary Fig.\ref{suppfig:branchlengths_scenario3}. Using the ML-tree leads to an overly liberal test; for very large internal branch lengths, 50\% of instances are deemed phylogenetic informative (Fig.\,\ref{fig:simulation}b), regardless of the sequence length. For external branches (Fig.\,\ref{fig:simulation}a), the empirical rejection rate is also too high at 15\%.

To roughly correct for the influence of the ML-inference, we can apply the Bonferroni correction to the number of pairs of leaves used to place branch $AB$. When the Bonferroni correction (see online methods \ref{msec:test}) is applied (ML-tree \& Bonferroni correction, {\color{MLtreeBonferroni} dark green}), the percentage of informative instances approaches the gold standard of scenario one. One word of caution: if many of the branch lengths of the tree are close to zero, then the test with Bonferroni correction becomes conservative and the rejection level falls below $5\%$ (Supplementary Fig.\,\ref{suppfig:simulation_shortBranches}).


Furthermore, if you examine the bootstrap support of the simulations for the 16-taxon tree in this scenario with an alignment length of 100bp, you will see that branch $AB$ was nearly always reconstructed (except for 8 instances out of 16000) and had a bootstrap support \citepMain{felsenstein1985confidence} of 100\% for all branch lengths greater than 0.5 time units, as shown in Supplementary Table \ref{supptab:bootstrap}. Typically, branches with high-bootstrap support are considered a real phylogenetic signal, justifying the separation between the two clusters in the tree. If you consult the results of the test for phylogenetic information, it becomes clear that bootstrap support cannot be used to justify the connection in the tree.

\medskip 

The results shown in Fig.\,\ref{fig:simulation} were generated under the assumption that no model misspecification occurs. Since the computation of the coherence is obviously depends on the selected model of sequence evolution (Eq.\,\eqref{eq:branch_likelihood_given_t}, \eqref{eq:branch_likelihood_at_site_t}), we  evolved sequences under a GTR-model {\color{olive}(REF)}  and then inferred an ML-tree  assuming either a Jukes-Cantor (JC) or Kimura two parameter (K2P) model. In both cases we observe a too liberal test (Supplementary Fig.\, \ref{suppfig:simulation_misspec}), with too many branches are called non-saturated, contrary to the gold-standard. However, when we replaced the JC or K2P model by the Fel81 model {\color{olive}(Felsenstein 1981)}, which incorporates the base composition of the alignment, the test performed well, and the rejection curves approach the gold-standard. 
Thus, the most critical aspect for the good performance of SatuTe is specifying the base composition, which is straightforward. The estimation of the substitution rate parameters in the rate matrix is less crucial.

From our limited set of simulations, we conclude that selecting a reasonable model of sequence evolution leads to a reasonable test behavior compared to the gold-standard. Moreover, the discovery of saturation at a branch should raise concern, as SatuTe tends to be too liberal.


    %%% Fig 2 - simulation %%%
	\begin{figure}[H]
            \centering
            \includegraphics[scale=0.5]{figure_simulated_data_main.pdf}
		     \caption{\textbf{Performance of SatuTe in simulations}. The plot shows the fraction of informative instances when testing the depicted branch in various scenarios: the true-tree ({\color{truetree}light red}), the true-tree topology with ML-estimated branch lengths ({\color{truetreeML} dark red}), and the inferred ML-tree including branch lengths ({\color{MLtree} light green}). The {\color{MLtreeBonferroni} dark green} curves represent the percentage of informative instances after applying the Bonferroni correction to the light green instances. In \textbf{a}, an external branch of a 5-taxon tree and in \textbf{b}, the central branch of a balanced 16-taxon tree is analyzed. Dashed, solid, and alternating dashed lines represent the results for 100, 1000, and 10000 sites, respectively.}
		     \label{fig:simulation}
        \end{figure}




\subsection*{SatuTe investigates the Tree of Life}

Finally, we aim to apply SatuTe to biological data to demonstrate its performance on real data and explore additional analysis possibilities. For this purpose, we are examining the work of \citetMain{Hug2016}, which tackles the complex problem of reconstructing the Tree of Life. 

\medskip

First, we consider the $3083$-taxon Tree of Life (ToL), which was inferred  using the LG plus a Gamma model of heterogeneity with 4 rate categories (LG+$\Gamma 4$ model). This ML-inference was based on a curated amino acid alignment of $2596$ sites, containing the concatenation of $16$ highly conserved proteins from both the large  (L) and small (S) ribosomal subunits (see Supplementary Table \ref{table:per_protein_summary}). For models with rate heterogeneity model, SatuTe performs a \textbf{per-category analysis} by splitting the alignment into the subalignments (in this case, four) and applying the test for phylogenetic information to the rescaled ToL. We specifically address the case of the inferred ML-tree with Bonferroni correction. Notably, all branches are phylogenetic informative across all categories, see Supplementary Section \ref{suppsec:2DT_per_category_analysis}.

\begin{figure}[ht]
            \centering
            \includegraphics[width=\linewidth]{figure_branch_saturation.pdf}
		     \caption{\textbf{Per-protein analysis: Saturation in the Tree of Life.} For each ribosomal protein, SatuTe tests for phylogenetic information on all branches of the pruned Tree of Life, which was reconstructed from the concatenated alignment of all $16$ ribosomal proteins. Branches identified as saturated for at least one of the proteins are highlighted in red.
        }
		     \label{fig:ToL_per_protein_analysis}
 \end{figure}

\medskip 

The alignment is composite, raising the question of where the selected ribosomal proteins are phylogenetic informative or not within the reconstructed tree topology. To address this, we can perform a \textbf{per-alignment-region analysis}. In this context, we consider the different proteins, though you can also use genes or other partitions for the alignment. To ensure comparable results for all proteins, we used the homogeneous LG model \citepMain{le2008}. For the analysis of each protein, we excluded taxa with no meaningful sequence data (i.e., sequences composed only of the characters “-” or “X”) and pruned the $3083$-taxon ToL accordingly. The number of excluded taxa ranged from 18 for L15 to 419 for S10, indicating a significant number of gaps. After ML-estimating the branch lengths in the given tree topology, we applied SatuTe, similar to the second scenario shown in the simulation (\textcolor{burgundy}{dark red}, Fig.\,\ref{fig:simulation}). In total, 715 branches are saturated for at least one of the 16 ribosomal proteins, highlighting the complexity of depicting branch saturation (Supplementary Table \ref{table:per_protein_summary}). As shown in Fig.\,\ref{fig:ToL_per_protein_analysis}, most saturated branches are in the Archaea subtree, which includes Eukaryota in this reconstruction. Remarkably, the four longest branches (namely the external branches of three Archaea Candidatus and the branch separating Bacteria from Archaea) are informative for all proteins. This may seem surprising considering that their branch lengths range from $0.92$ to $1.22$ substitutions per site, while the shortest protein has only $74$ sites. For more details of this analysis see Supplementary Section \ref{suppsec:2DT_per_region_analysis}.

\begin{figure}[ht]
            \centering
            \includegraphics[width=\linewidth]{Figure_2DT_sliding_window_2024_06_21.pdf}
		     \caption{\textbf{Per-category and Sliding-window Analysis of Specific Branches:} \textbf{a}, 2D tree with $3083$ taxa obtained from the ribosomal protein alignment. The branches leading to Eukaryota and Yeast were analysed. \textbf{b}, For each of the four rate categories, SatuTe tested the branches for phylogenetic information using the subalignments and rescaled tree. The per-category z-scores are marked as dashed lines, and the number of sites per category is represented by bars. \textbf{c}, Plot of z-scores from $36$bp sliding-window analysis along the protein alignment. 
}  \label{fig:EUK} 
\end{figure}

\medskip 

SatuTe also enables \textbf{branch-specific analyses}. In the following, we focus on the branch leading to the Eukaryota, as the placement of the Eukaryota in the ToL is subject of ongoing debate. The reconstruction of the ToL depicted in Fig.\,\ref{fig:ToL_per_protein_analysis} positions Eukaryota as a subdomain of Archaea, thereby splitting all taxa into two domains (2D): Bacteria and Archaea. The initial phylogenetic analyses based on the 16S rRNA gene \citepMain{woese} classified Eukaryota as a monophyletic domain separate from Bacteria and Archaea, thus dividing all taxa into three domains (3D). We used SatuTe to determine whether this disagreement is due to saturation. To investigate this, we return to our initial (per-category) analysis and examine the results for the branch leading to Eukaryota in the 2D tree (Fig.\,\ref{fig:EUK}a). This analysis shows that under the 2D tree, no rate category is affected by saturation (Fig.\,\ref{fig:EUK}b). To gain a more detailed understanding of the changes in phylogenetic information, we additionally performed a \textbf{sliding-window analysis} with a window size of $36$bp. This analysis detects only a minority of sites affected by saturation, primarily at the extremes of the protein sequences (Fig.\,\ref{fig:EUK}c). We  also conducted the same analysis for an external branch, namely the Yeast branch, and found very similar results.

\medskip

The sliding-window approach can also be applied to individual proteins, using the obtained z-scores to compare different branches and identify potential information loss by examining the differences (see Supplementary Information\,\ref{suppsec:2DT_branch_zscore_diffs} for per-protein z-score differences between Eukaryota and Yeast). %Supplementary Fig.\,\ref{fig:Eukaryota_Yeast}  illustrates that protein L16 experiences the most significant information loss at the branch leading to Eukaryota, based on the per-protein \textbf{z-score difference} between Eukaryota and Yeast.
But in addition to the branch comparison, \textbf{z-score differences} allow us to determine whether each protein supports the 2D or 3D topology of the ToL. We curated the 3D-rearranged tree (Fig.\,\ref{fig:z-score_differences}), placing Eukaryota as a separate third domain distinct from Bacteria and Archaea. As shown in Fig.\,\ref{fig:z-score_differences}, for the branch leading to Eukaryota,  three out of the sixteen protein (L15, L22, and S10) favor the 3D-rearranged tree. Of the eleven proteins from the large ribosomal subunit, six (L2, L3, L5, L6, L16, and L18) weakly favor the 2D tree, with z-score differences ranging from 0.023 for protein L2 to 0.106 for protein L18.

The support for the 3D-rearranged tree may be partly due to gaps, as the z-score difference in Fig.\,\ref{fig:z-score_differences} shows a negative correlation coefficient of $-0.60$ with the fraction of gaps for each protein in the Asgard subalignment. A dataset with more and better-sampled Asgard sequences or a higher site selection stringency as discussed in \citepMain{Hug2016} might favor the 2D tree for all proteins considered here. Notably, proteins of the small ribosomal subunit more robustly support the 2D tree, with all differences greater than a $0.25$ z-score except for S10.

\begin{figure}[ht]
            \centering
            \includegraphics[width=\linewidth]{figure_zscore_differences_2024_06_21.pdf}
		     \caption{\textbf{Per-protein z-scores comparison to a rearranged topology.} Scatter plot showing the fraction of gapped sites in the Asgard superphylum and the z-score difference for the branch leading to Eukaryota between the 2D tree (Fig. \ref{fig:EUK}a) and 3D-rearranged tree (bottom left) for each ribosomal protein. The Asgard alignment includes 7 sequences, representing Lokiarchaeota and Thorarchaeota. The 3D-rearranged tree obtained from the 2D tree by replacing the black dashed branch with the green dashed branch.  The Pearson correlation coefficient is $\rho = -0.60$.
}  \label{fig:z-score_differences} 
\end{figure}
\medskip

 
As previously mentioned,  while 2D trees are consistently inferred from ribosomal protein alignments \citepMain{williams2D}, some 3D trees still emerge depending on the chosen alignment \citepMain{Hug2016, gribaldo2010}. \citetMain{Hug2016} also provided a $1871$-taxon 3D tree reconstructed from a 16S rRNA gene alignment of $1947$ sites under GTR+$\Gamma 4$ model (see Supplementary Information \ref{suppsec:rRNA_based_3DT} for similar analyses). For the fastest evolving rate category (rate $2.73$), $107$ branches are classified as saturated. Among these, $13$ are internal branches, including the branch leading to Eukaryota. The sliding-window analysis for this branch leading to Eukaryota shows that the fraction of saturated $36$bp windows is $43\%$ ($822$ out of $1912$). However, since still more than half of the 16S rRNA gene alignment is informative for this branch, the 3D topology is not solely an artifact of saturation and should not be excluded on this basis alone. Paradoxically, while the 16S rRNA belongs to the small ribosomal subunit, the 16S rRNA gene favours the 3D ToL for $6$ out of $9$ rRNA regions when an analogous 2D-rearranged tree is constructed (Fig.\,\ref{suppfig:rRNA_2D_3D}). This suggests that explaining the disagreement between the 16S rRNA and its surrounding proteins (Fig.\,\ref{fig:z-score_differences}) could be a useful hint to find a reliable reconstruction of the ToL.
    
\bigskip









\section*{Discussion}


%\begin{itemize}
    %\item simulation and real data
    %\item Considering these z-score differences, the biggest alignment we could construct favouring the 3D hypothesis in log-likelihood contained $5$ proteins (the 3D-supporting L15, L22, S10 and the weakly  2D-supporting L2, L4, see Supplementary Table \ref{table:1}).
    %\item  A dataset with more and better-sampled Asgard sequences or a higher site selection stringency as discussed in \citepMain{Hug2016} might favor the 2D tree for all proteins considered here.
%\end{itemize}


Since SatuTe can measure the phylogenetic information for any subalignment and any branch in a phylogeny, the number of possible analyses and graphical representations is immense. We demonstrated the use of SatuTe with both simulated and real data. The z-scores provided by SatuTe are straightforward and easy to interpret, facilitating comparisons between different alignments, branches and candidate topologies, as well as visualizing regional changes in information content. We encourage practitioners to propose new analyses, as SatuTe provides the necessary elements of the test statistic for any given branch. SatuTe is a complementary tool that can be combined with established reliability assessments, such as the bootstrap by \citepMain{felsenstein1985confidence}, which may give high bootstrap support values to saturated branches (Supplementary Table \ref{supptab:bootstrap}).


%In this work, we have left aside the problem of tree reconstruction, and focused on measuring the information gluing together the subtrees of a given tree. In this sense, we emphasise that SatuTe does not infer a tree, but rather evidentiates the most solid part of a phylogeny, in a complementary analysis to e.g. the phylogenetic bootstrap \citepMain{felsenstein1985confidence}. As we have shown, the z-scores provided by SatuTe are simple and comprehensible, allowing to compare different alignments, branches of a given tree, and candidate topologies. 


\medskip

Regarding the 2D Tree of Life obtained from the ribosomal protein alignment, SatuTe showed that it is mostly unaffected by saturation.
%, for the protein with the highest fraction of saturated branches was L16, with only $3.2\%$ ($192/5994$).
For the placement of Eukaryota in the ToL, only a minority of proteins preferred the 3D-rearranged tree. Although resolving the 2D versus 3D conflict is beyond our scope, the fact that ribosomal proteins can be selected to favour either topology highlights the risk of using arbitrary subsets of proteins or genes as definitive evidence for the true ToL. Our results suggest that the 2D tree could receive stronger support with the inclusion of more and better-sampled Asgard sequences in the alignment or with stricter site selection criteria. Even when gaps are treated as in the maximum likelihood calculations, they still influence the test.
 %\citepMain{williams2D,cunha3D} 

\medskip

Interestingly, we attributed the near absence of saturation in the ribosomal protein alignment to largest non-zero eigenvalue of the LG model, $\lambda_1 \approx -0.27$,  which results in a slow exponential decay $e^{\lambda_1 t}$. We think the advantage of amino acid models over nucleotide models for deep phylogenies lies precisely in these larger eigenvalues, a point not discussed by \citepMain{kapli2023}. Focusing on the eigenvalues of the rate matrix, well known among theoreticians \citepMain{levin2017markov,elchanan2003}, offers a valuable perspective for discussing deep phylogenies.

\medskip



%Apart from the SatuTe implementation, our measure of information yields a simple condition for the convergence of the Newton method when ML-estimating branch lengths (Corollary \ref{prop:decission_tree}). Future research should formalise the relationship between the SatuTe z-score and the MLE tree reconstruction.

%{\color{truetree}The bootstrap or its derivatives \citepMain{lemoine2018, hoang2017, Stamatakis2008} are  not an adequate tool to detect branch saturation either, since intuitively the split induced by a long branch should be easily reproducible and have a high bootstrap support. Differently said, the bootstrap yields a difficult diagnosis: The overall bootstrap values may decrease due to the misplacement of the saturated branch, while the bootstrap value of the split induced by the saturated branch itself should be high.}{\color{blue} This is rather a result, rather to the discussion. Only bootstrap.}

%{\color{blue} Insist in the discussion that as long as reconstruction is reasonably well specified, saturated branches detected by SatuTe are indeed source of concern and should be treated with caution. }

Our measure of information can be generalised for non-stationary and non-reversible models using the $\pi$-scalar product between the likelihood vectors, as described by \citepMain{manuel2022}. However, dropping those assumptions leads to root identifiability. This implies that the test for phylogenetic information does not depend only on the parent node of each subtree, but also on the placement of the root. While this complication makes the statement of the test slightly more technical, its implementation remains feasible for future research. 



\section*{Online content}

{\color{blue}Any methods, additional references, Nature Portfolio reporting summaries, source data, extended data, supplementary information, acknowledgements, peer review information; details of author contributions and competing interests; and statements of data and code availability are available at  }


\backmatter
 
% To change the title from References to Bibliography:


%\renewcommand\refname{Bibliography}
\bibliographystyleMain{apalike}
\bibliographyMain{biblio.bib} 


%\begin{refsection}
%\nocite{manuel2022}
%\printbibliography
%\end{refsection}